% CANDIDATE computational note. Not a theorem. Does not edit main.tex.
% XeLaTeX. Sources cited below were opened on 2026-10-04.
\documentclass[11pt]{article}
\usepackage{fontspec}
\usepackage{amsmath,amssymb}
\usepackage[margin=1in]{geometry}
\usepackage{booktabs}
\usepackage{hyperref}
\setmainfont{Latin Modern Roman}
\hypersetup{colorlinks=true,urlcolor=blue,linkcolor=blue}

\title{COMPUTATIONAL: polynomial-calculus space and RAM-versus-BP width\\
\large PC\_RAM1, not a theorem}
\author{Grok Bot, identity \texttt{j15marti}\\
j15marti/pnp, branch \texttt{main}}
\date{4 October 2026, about 11:57 AM PT}

\begin{document}
\maketitle

P versus NP remains open.

This note runs the two experiments left open by
\texttt{research/agents/grok/2026-10-04\_1149\_GROK\_PROOFSPACE1.tex}
(commit \texttt{ddccc9827409d3c183fa262015340a4a4db47851},
blob \texttt{22ed91add377f55ff3837a64e15011274e76ac18}):
polynomial-calculus space in the literature, and a finite
TEST-PROOFSPACE-D48 / RAM-versus-branching-program check on one
frozen \(D(4,8)\) CNF together with a recomputed ordered DAG witness.
It does not edit root \texttt{main.tex}. It does not add a workflow.
It does not claim that P versus NP is solved. No number below is
invented: every COMPUTATIONAL figure was produced by a process run in
this pass, and every literature claim is taken from a source opened
in this pass.

\section*{Standing project facts, not revived}
\begin{itemize}
\item P versus NP remains open.
\item Product-span rank and the compiler into a charged \((T,S,R)\) model remain \textbf{OPEN}.
\item Semantic cut \(\ge 48\) on negative-star \(D(4,8)\) is \textbf{RETIRED} (exact cut \(8\)); contact count \(48\) stays exact inside that architecture.
\item \(Q_4\) as a read-2 lower-bound candidate is \textbf{RETIRED}.
\item Contacts agree with branching programs up to a polynomial and still do not close the compiler gap.
\end{itemize}

\section*{Part A. Literature actually opened}

\subsection*{Polynomial-calculus space (degree, monomial space, total space)}
Opened and read:
\begin{itemize}
\item Nicola Galesi, Leszek~A.\ Ko\l{}odziejczyk, and Neil Thapen,
\emph{Polynomial Calculus Space and Resolution Width},
\emph{Theory of Computing} 21(6), 2025, pages 1--29.
PDF: \url{https://theoryofcomputing.org/articles/v021a006/v021a006.pdf}.
\item Massimo Lauria, Mladen Mik\v{s}a, Jakob Nordstr\"om, and Marc Vinyals,
\emph{A Simplified Proof for the Monomial Space Complexity of Random 4-CNF Formulas in Polynomial Calculus},
\emph{Theory of Computing}, dated 29 December 2025 in the opened PDF.
PDF: \url{https://jakobnordstrom.se/docs/publications/SimplifiedProofPCRspaceRandom4CNF_ToC.pdf}.
\end{itemize}

What is proved, in one sentence each.
Galesi--Ko\l{}odziejczyk--Thapen, Theorem~6.4 as read: if a \(k\)-CNF has a PCR refutation of term space \(s\) over some field, then it has a resolution refutation of width \(s^2-s+k\); equivalently, resolution width \(w\) forces PCR monomial/term space at least on the order of \(\sqrt{w}\).
Lauria--Mik\v{s}a--Nordstr\"om--Vinyals, Corollary~3.4 as read: a random \(k\)-CNF with \(k\ge 4\) and fixed density requires \(\Omega(n)\) monomial space in PCR asymptotically almost surely.
The same opened ToC paper of Galesi--Ko\l{}odziejczyk--Thapen also records Bonacina's total-space bound for resolution (total space at least on the order of width squared) with a simplified proof.

What those theorems do \textbf{not} say about arbitrary SAT.
They are statements about PCR (or resolution) configurations on CNF refutations.
They do not put SAT outside P.
They do not identify PCR monomial space with the project's product-span obstruction.
They do not prove that PCR degree is a lower bound on monomial space: both opened papers explicitly leave open whether the Atserias--Dalmau width-versus-space relation has a degree-versus-monomial-space analogue without a square-root loss.
``Width approaches degree'' is therefore still a candidate meeting line, not a theorem of the sources opened here.
No connection from PCR space to \(D(n,m)\) product-span was read.

\subsection*{RAM space versus logarithm of branching-program width}
Opened and read:
\begin{itemize}
\item Paul Beame, Michael Saks, and Jayram~S.\ Thathachar,
\emph{Time-Space Tradeoffs for Branching Programs}.
PDF: \url{https://homes.cs.washington.edu/~beame/papers/focsbranch.pdf}.
\end{itemize}

What is proved (or defined), in one sentence.
In the opened Beame--Saks--Thathachar manuscript, branching-program time is path length and branching-program space is defined as the logarithm of the number of nodes; that is a model definition inside branching programs, together with concrete length-versus-size lower bounds for specific functions (including length \(1.0178n\) forcing size \(\ge 2^{n/3}\) for \(BQF\) on a Sylvester matrix over \(\mathrm{GF}(3)\)).

What that does \textbf{not} say about arbitrary SAT.
It does not prove that the work-space of a uniform random-access Turing machine equals the logarithm of the width (or size) of a branching program for SAT.
Fortnow-style RAM time-space lower bounds for SAT, opened in the previous PROOFSPACE1 note, live in a different model; the opened Beame--Saks--Thathachar text itself says Fortnow's bound does not transfer to nonuniform branching programs.
So ``RAM space \(= \log(W_{\mathrm{BP}})\)'' remains an analogy / candidate boundary, not a verified simulation theorem for SAT.
Status of that equality as a general claim: \textbf{not a theorem in the sources opened}; on the finite \(D(4,8)\) witness below it is numerically false for the direct mask-and-count charge.

\section*{Part B. Computation}

\subsection*{Repository semantics used}
Read the PROOFSPACE1 note and the repository encoding of \(D(4,8)\):
\texttt{experiments/universal-dag/d5\_m8\_exact.cpp}
(blob \texttt{1ca7c97ca45fce506411d246f0f4c499a1352de1} at the parent commit)
and the Python reconstruction
\texttt{ord\_dag.py} matching that C++ accounting.
Input bits are presence of the \(32\) simple signed \(3\)-clauses on \(n=4\) variables; promise weight \(m=8\); output \(1\) means the selected clauses falsify every assignment.
Population \(\binom{32}{8}=10518300\).
The full population was not treated as a universal certificate.
One named frozen formula was used for the proof-system coordinates; the ordered DAG was recomputed for the RAM-versus-width coordinates.

\subsection*{Named frozen CNF (TEST-PROOFSPACE-D48)}
\textbf{Name.} \(F_\star\): the weight-\(8\) selection whose clause indices in the sign-round-robin order are
\[
(0,4,8,12,16,20,24,28).
\]
Explicit clauses (variables \(\{1,2,3,4\}\); this selection is all eight sign patterns on \(\{2,3,4\}\), omitting variable \(1\)):
\begin{align*}
&(x_2\lor x_3\lor x_4),\ 
(\lnot x_2\lor x_3\lor x_4),\ 
(x_2\lor \lnot x_3\lor x_4),\ 
(\lnot x_2\lor \lnot x_3\lor x_4),\\
&(x_2\lor x_3\lor \lnot x_4),\ 
(\lnot x_2\lor x_3\lor \lnot x_4),\ 
(x_2\lor \lnot x_3\lor \lnot x_4),\ 
(\lnot x_2\lor \lnot x_3\lor \lnot x_4).
\end{align*}
Size: \(4\) variables, \(8\) clauses, a \(3\)-CNF.
Machine-checked unsatisfiable by exhaustive search of all \(16\) assignments, and by mask-union equal to \(\mathtt{0xffff}\).
This is one of the \(272\) UNSAT weight-\(8\) formulas in the \(D(4,8)\) accounting; it is not a certificate for all \(10518300\) inputs.

\textbf{Resolution width (COMPUTATIONAL).}
Minimum resolution width of \(F_\star\) is \(3\).
Reason: every axiom has width \(3\), so width is at least \(3\); a width-\(3\) closure derives the empty clause (closure size \(27\) clauses in the saturating implementation run here).
Atserias--Dalmau lower bound on clause space: \(\mathrm{CSpace}\ge w_{\mathrm{Res}}-3=0\), which is vacuous on this instance.

\textbf{Characteristic-2 PC degree (COMPUTATIONAL).}
Minimum polynomial-calculus degree of \(F_\star\) over \(\mathrm{GF}(2)\) (Boolean ring, \(x_i^2=x_i\)) is \(3\).
Reason: each axiom translates to a degree-\(3\) monomial equation, so degree is at least \(3\); a degree-\(3\) Buchberger-style closure over \(\mathrm{GF}(2)\) derives \(1\).
On this one frozen formula, resolution width and characteristic-2 PC degree meet at \(3\).
That meeting is a fact about \(F_\star\) only.
It is not a general theorem that width approaches degree, and it does not touch product-span.

\subsection*{Recomputed ordered DAG and RAM-style charge}
Sign-round-robin reduced ordered clause-query DAG for \(n=4\), \(m=8\), same accounting as the C++ / ITER1 witness, recomputed in this pass:
\begin{align*}
\text{nonterminal nodes}&=720,\\
\text{peak level width}&=66
\quad\text{(levels \(15\) and \(16\) both have width \(66\))},\\
\log_2 66 &= 6.044394119358453.
\end{align*}
At level \(16\), live selected-counts are \(k\in\{2,3,4,5,6\}\) with sizes \(1+16+32+16+1=66\).
Cell \((t,k)=(16,4)\): \(K=32\) decision identities.
Continuation-incompatibility certificate recomputed: all \(\binom{32}{2}=496\) pairs are separated by some common weight-\(4\) continuation of the remaining \(16\) clauses.
So \(\log_2 K=5\).
All of these numbers are \textbf{COMPUTATIONAL} (finite, machine-checked this pass). They match the ITER1 committed witness and are not trusted unread.

RAM-style mutable space charge for the direct decider that stores the coverage mask and the selected count (excluding the program counter):
\[
S_{\mathrm{mask}}=16\text{ bits},\qquad
S_k=\lceil\log_2 9\rceil=4\text{ bits},\qquad
S_{\mathrm{mutable}}=20\text{ bits}.
\]
Including a clause-index program counter \(\lceil\log_2 33\rceil=6\) gives \(26\) bits.
Comparison:
\[
S_{\mathrm{mutable}}=20 \;\neq\; \log_2 66\approx 6.044,\qquad
S_{\mathrm{mutable}}=20 \;\neq\; \log_2 K=5.
\]
So on this finite ordered \(D(4,8)\) witness, the direct RAM-style space charge is \textbf{not} equal to the logarithm of branching-program level width, nor to \(\log_2 K\).
The experiment therefore does \textbf{not} connect RAM space to \(\log(W_{\mathrm{BP}})\) as an equality; it exhibits a concrete numerical disagreement under the stated charge.
It also does not connect PCR space to the product-span obstruction: no product-span matrix was formed, and no PCR-space number was claimed for the whole \(D(4,8)\) architecture.

\section*{Status ledger}
\begin{center}
\begin{tabular}{@{}llp{2.4in}@{}}
\toprule
Item & Status & Note \\
\midrule
P versus NP & OPEN & opening sentence \\
PCR space \(\sqrt{w}\) vs resolution width & THEOREM & Galesi--Ko\l{}odziejczyk--Thapen, opened \\
random \(4\)-CNF PCR monomial space \(\Omega(n)\) & THEOREM & Lauria--Mik\v{s}a--Nordstr\"om--Vinyals, opened \\
degree lower-bounds PCR monomial space & OPEN & both opened PC-space papers \\
width approaches degree (general) & OPEN / CANDIDATE & not a theorem opened here \\
RAM space \(=\log(W_{\mathrm{BP}})\) for SAT & not a theorem & BP definition only; RAM is separate \\
\(F_\star\) resolution width \(=3\) & COMPUTATIONAL & one frozen CNF \\
\(F_\star\) char-2 PC degree \(=3\) & COMPUTATIONAL & one frozen CNF \\
DAG width \(66\), \(K=32\), \(\log_2 66\), \(\log_2 K=5\) & COMPUTATIONAL & recomputed \\
RAM mutable charge \(20\) vs \(\log_2 66\) & COMPUTATIONAL & not equal \\
PC space \(\leftrightarrow\) product-span & not connected & this pass \\
RAM space \(\leftrightarrow\log W_{\mathrm{BP}}\) & not connected as equality & this pass \\
general \((T,S,R)\) tradeoff & OPEN & \\
compiler into charged model & OPEN & \\
TEST-PROOFSPACE-D48 on \(F_\star\) & COMPUTATIONAL & was UNRUN in PROOFSPACE1 \\
\bottomrule
\end{tabular}
\end{center}

\section*{What this note does not say}
It does not solve P versus NP.
It does not promote \(F_\star\) into a universal certificate for the \(10518300\)-instance promise problem.
It does not claim that resolution width equals PC degree in general.
It does not claim that RAM work space equals the log of branching-program width in general.
It does not claim that PCR space bounds the project's product-span rank.
The two load-bearing arrows---product-span and the compiler---remain OPEN.

\section*{Plain-language summary}
P versus NP is still open.
We opened real PCR-space papers: they relate monomial space to resolution width with a square-root loss, and they leave open whether degree itself lower-bounds space; none of that talks about our product-span obstruction or puts SAT outside P.
We opened the Beame--Saks--Thathachar branching-program paper: there ``space'' means log of the number of nodes inside the branching-program model, which is not a proved equality between RAM Turing-machine space and branching-program width for SAT.
On one small frozen unsatisfiable \(8\)-clause formula taken from \(D(4,8)\), resolution width and characteristic-2 PC degree both came out to \(3\) (computational, that formula only).
On the recomputed ordered \(D(4,8)\) DAG we again saw level width \(66\) and a tight cell of size \(32\); a direct RAM-style mask-and-count charge used \(20\) bits, which is not equal to \(\log_2 66\).
So neither experiment connected RAM space to log of BP width as an equality, and neither connected PC space to our product-span obstruction.
The general tradeoff and the compiler stay open.

% END_OF_VERIFIED_REPORT::GROK::2026-10-04_1157_PC_RAM1
\end{document}
